\documentclass[twoside]{article}
\setlength{\oddsidemargin}{0.25 in}
\setlength{\evensidemargin}{-0.25 in}
\setlength{\topmargin}{-0.6 in}
\setlength{\textwidth}{6.5 in}
\setlength{\textheight}{8.5 in}
\setlength{\headsep}{0.75 in}
\setlength{\parindent}{0 in}
\setlength{\parskip}{0.1 in}
\usepackage[spanish]{babel}
\selectlanguage{spanish}
\usepackage[utf8]{inputenc}


\usepackage{amsmath,amsfonts,graphicx}


\newcounter{lecnum}
\renewcommand{\thepage}{\thelecnum-\arabic{page}}
\renewcommand{\thesection}{\thelecnum.\arabic{section}}
\renewcommand{\theequation}{\thelecnum.\arabic{equation}}
\renewcommand{\thefigure}{\thelecnum.\arabic{figure}}
\renewcommand{\thetable}{\thelecnum.\arabic{table}}

\newcommand{\lecture}[4]{
   \pagestyle{myheadings}
   \thispagestyle{plain}
   \newpage
   \setcounter{lecnum}{#1}
   \setcounter{page}{1}
   \noindent
   \begin{center}
   \framebox{
      \vbox{\vspace{2mm}
    \hbox to 6.28in { {\bf 1076: An\'alisis Matem\'atica 1 
		\hfill Tercer Semestre 2015} }
       \vspace{4mm}
       \hbox to 6.28in { {\Large \hfill Lección  #1: #2  \hfill} }
       \vspace{2mm}
       \hbox to 6.28in { {\it Profesor: #3 \hfill  Estudiante: #4} }
      \vspace{2mm}}
   }
   \end{center}
   \markboth{Lecture #1: #2}{Lecture #1: #2}

   {\bf Responsabilidad}: {\it  Estas notas no han sido sometidos al escrutinio de costumbre reservado para publicaciones formales. Éstas NO  pueden ser distribuidas fuera de esta clase.}

   \vspace*{4mm}
}

\renewcommand{\cite}[1]{[#1]}
\def\beginrefs{\begin{list}%
        {[\arabic{equation}]}{\usecounter{equation}
         \setlength{\leftmargin}{2.0truecm}\setlength{\labelsep}{0.4truecm}%
         \setlength{\labelwidth}{1.6truecm}}}
\def\endrefs{\end{list}}
\def\bibentry#1{\item[\hbox{[#1]}]}


\usepackage{amsthm}
 
\theoremstyle{theorem}
\newtheorem{theorem}{Teorema}[lecnum]
\newtheorem{lemma}[theorem]{Lema}
\newtheorem{proposition}[theorem]{Proposición}
\newtheorem{corollary}[theorem]{Corolário}

\theoremstyle{definition}
\newtheorem{definition}{Definición}[theorem]
\newtheorem{exercise}{Ejercício}[theorem]
 
\theoremstyle{remark}
\newtheorem{remark}{Observación}[theorem]



\newcommand\E{\mathbb{E}}

\begin{document}

%FILL IN THE RIGHT INFO.
%\lecture{**LECTURE-NUMBER**}{**DATE**}{**LECTURER**}{**SCRIBE**}
\lecture{1}{August 24}{Humberto Rafeiro}{scribe-name  }
%\footnotetext{These notes are partially based on those of Nigel Mansell.}

% **** YOUR NOTES GO HERE:


\begin{abstract}
Adicionar un pequeño resumen, entre dos a tres líneas.
\end{abstract}

\tableofcontents

$\uplus, \cup$

\section{Some theorems and stuff} % Don't be this informal in your notes!

We now delve right into the proof.

\begin{lemma}
This is the first lemma of the lecture.
\end{lemma}

\begin{proof}
The proof is by induction on $\ldots$.
\end{proof}

\subsection{A few items of note}

Here is an itemized list:
\begin{itemize}
\item this is the first item;
\item this is the second item.
\end{itemize}

Here is an enumerated list:
\begin{enumerate}
\item this is the first item;
\item this is the second item.
\end{enumerate}

Here is an exercise:

\begin{exercise}
Show that ${\rm P}\ne{\rm NP}$.

Here is how to define things in the proper mathematical style.
Let $f_k$ be the $AND-OR$ function, defined by

\[ f_k(x_1, x_2, \ldots, x_{2^k}) = \left\{ \begin{array}{ll}

	x_1 & \mbox{if $k = 0$;} \\

	AND(f_{k-1}(x_1, \ldots, x_{2^{k-1}}),
	   f_{k-1}(x_{2^{k-1} + 1}, \ldots, x_{2^k}))
	 & \mbox{if $k$ is even;} \\

	OR(f_{k-1}(x_1, \ldots, x_{2^{k-1}}),
	   f_{k-1}(x_{2^{k-1} + 1}, \ldots, x_{2^k}))	
	& \mbox{otherwise.} 
	\end{array}
	\right. \]
\end{exercise}

\begin{theorem}
This is the first theorem.
\end{theorem}

\begin{proof}
This is the proof of the first theorem. We show how to write pseudo-code now.
%*** USE PSEUDO-CODE ONLY IF IT IS CLEARER THAN AN ENGLISH DESCRIPTION

Consider a comparison between $x$ and~$y$:
\begin{tabbing}
\hspace*{.25in} \= \hspace*{.25in} \= \hspace*{.25in} \= \hspace*{.25in} \= \hspace*{.25in} \=\kill
\>{\bf if} $x$ or $y$ or both are in $S$ {\bf then } \\
\>\> answer accordingly \\
\>{\bf else} \\
\>\>    Make the element with the larger score (say $x$) win the comparison \\
\>\> {\bf if} $F(x) + F(y) < \frac{n}{t-1}$ {\bf then} \\%
\>\>\> $F(x) \leftarrow F(x) + F(y)$ \\
\>\>\> $F(y) \leftarrow 0$ \\
\>\> {\bf else}  \\
\>\>\> $S \leftarrow S \cup \{ x \} $ \\
\>\>\> $r \leftarrow r+1$ \\
\>\> {\bf endif} \\
\>{\bf endif} 
\end{tabbing}

This concludes the proof.
\end{proof}


\section{Next topic}

Here is a citation, just for fun~\cite{CW87}.

\section*{Bibliografía}
\beginrefs
\bibentry{CW87}{\sc D.~Coppersmith} and {\sc S.~Winograd}, 
``Matrix multiplication via arithmetic progressions,''
{\it Proceedings of the 19th ACM Symposium on Theory of Computing},
1987, pp.~1--6.
\endrefs

% **** THIS ENDS THE EXAMPLES. DON'T DELETE THE FOLLOWING LINE:

\end{document}





